\chapter{Overview}

\begin{center}
	\includegraphics[width=0.5\textwidth]{Images/aberfitness.png}
\end{center}

\textit{Aber Fitness} is a web application developed using Microsoft's \textit{.NET Core} and Oracle's \textit{Java Enterprise Edition} (henceforth referred to as Java EE). The project aims to provide a service to encourage fitness and promote engagement with sporting activities amongst the users of the application, offering functionality such as graphing fitness data gathered by owners of \textit{Fitbit} devices, the ability to challenge other users to competitions and a sport ladder system with tight integration into a bespoke facility booking system. \textit{Aber Fitness} aims to offer everything that would be needed by a sporty and active person in order to bring their sporting activities into a digital platform and also to enhance their use of devices they already own, such as \textit{Fitbit} devices or smart watches such as the \textit{Apple Watch}.

At launch the system will ingest activity data automatically from \textit{Fitbit}, with the capability of easily implementing other health data provider services at a later date due to the modular nature of the data ingest system. Once normalised this activity data will be used throughout the various subsystems of \textit{Aber Fitness}, providing users with functionality such as a dashboard overview of their activity over the last hour, day, week, etc. as well as integrating tightly into the challenges system to add a competitive aspect to the system in to keep users engaged with both the platform itself and keeping fit in general. 

Aber Fitness will also feature a number of supporting tools, such as: a user management system with two-factor authentication, the ability for users to view a comprehensive audit trail of when and how their data was used and accessed, and a global layout template with a dynamic sidebar which can be adjusted by administrators and changes instantly propagated to other microservices. 

The platform is also designed to be scalable and easily deployed through the use of \textit{Docker} and the segregation of each aspect of the system into individual microservices. 

\textbf{TODO: Possibly needs padding out a bit more.}